\documentclass[11pt]{article}
\usepackage[T1]{fontenc}
\usepackage{lmodern}
\usepackage[margin=1.05in]{geometry}
\usepackage{amsmath,amssymb,amsthm,mathtools}
\usepackage{microtype}
\usepackage{booktabs,array}
\usepackage[colorlinks=true,linkcolor=black,citecolor=black,urlcolor=blue]{hyperref}
\newtheorem{theorem}{Theorem}
\newtheorem{lemma}[theorem]{Lemma}
\newtheorem{corollary}[theorem]{Corollary}
\theoremstyle{remark}
\newtheorem{remark}[theorem]{Remark}
\DeclareMathOperator{\ord}{ord}
\newcommand{\vp}{v_P}
\newcommand{\phithree}[1]{(#1^2+#1+1)}
\newcommand{\phisix}[1]{(#1^2-#1+1)}
\title{Orders Three and Six at the Largest Prime Divisor\protect\\ of an Odd Perfect Number}
\author{}
\date{October 9, 2026}
\begin{document}
\maketitle
\vspace{-1.3em}
\begin{abstract}
Let $N$ be an odd perfect number and $P$ its largest prime divisor. We prove that components whose bases have multiplicative order $3$ or $6$ modulo $P$ contribute at most two to $v_P(\sigma(N))$, and at most one if $3\nmid N$. We sharpen this bound using the available $3$-adic valuation and the Euler prime. A further refinement concerns distinct fourth-power prime divisors: if $P\equiv1\pmod3$, then $K+v_3(N)\ge2v_3(a_P+1)+\delta$, with $K$ and $\delta$ defined explicitly below. In the opposite direction, cyclotomic values give an upper bound for the contribution outside orders $3$ and $6$. If there are $M$ contributing components and $m$ is the odd part of $P-1$, their total contribution is at most $M(\varphi(m)-1)$. All stated bounds have accompanying compiled Lean 4 proofs. Publication priority is not claimed.
\end{abstract}

\section{Statement and notation}

Write $\sigma(n)=\sum_{d\mid n}d$. Throughout, $N$ denotes an odd perfect number, so that $\sigma(N)=2N$, and
\[
 N=\prod_{q\in S}q^{a_q},\qquad
 S=\{q:q\text{ is prime and }q\mid N\},\qquad P=\max S.
\]
For a prime $r$, the valuation $v_r(m)$ is the exponent of $r$ in the positive integer $m$. For $q\in S\setminus\{P\}$, set
\[
 d_q=\ord_P(q),\qquad w_q=v_P\bigl(\sigma(q^{a_q})\bigr),
 \qquad
 H=\{q\in S:q<P,\ d_q\notin\{3,6\}\}.
\]
A component $q^{a_q}$ contributes to $P$ when $w_q>0$. Define
\[
 c(N)=\begin{cases}2,&3\mid N,\\1,&3\nmid N.\end{cases}
\]

\begin{theorem}\label{thm:main}
The contributions to $P$ satisfy
\begin{equation}\label{eq:lowbound}
 \sum_{\substack{q\in S,\ q<P\\d_q\in\{3,6\}}}w_q\le c(N).
\end{equation}
Consequently,
\begin{equation}\label{eq:main}
 \boxed{\displaystyle\sum_{q\in H}w_q\ge a_P-c(N).}
\end{equation}
\end{theorem}

The valuation formulas underlying factor-chain arguments are classical; see, for example, Voight~\cite{Voight}. Relations between distinct bases sharing a factor of $x^2+x+1$ occur in Clayton and Hansen~\cite{ClaytonHansen}. The assertion here is the combined bound for arbitrary component exponents, including contributions of order $6$. No claim of publication priority is made.

\section{Valuations at the largest prime}

\begin{lemma}\label{lem:noexponent}
For every $q\in S$, one has $P\nmid a_q+1$.
\end{lemma}
\begin{proof}
Suppose that $P\mid a_q+1$. Then
\[
 C=1+q+\cdots+q^{P-1}
\]
divides $\sigma(q^{a_q})$. The integer $C$ is odd and greater than one. Since $q\le P$ is prime, $q\not\equiv1\pmod P$. Fermat's little theorem, including the case $q=P$, gives $C\equiv1\pmod P$.

Let $r$ be a prime divisor of $C$. Then $r$ is odd and $r\ne P$. Also $r\nmid q$. If $q\equiv1\pmod r$, then $C\equiv P\pmod r$, contrary to $r\ne P$. Thus $q^P\equiv1\pmod r$ and $q\not\equiv1\pmod r$. As $P$ is prime, $\ord_r(q)=P$. It follows that $P\mid r-1$, and hence $r>P$.

On the other hand, $r\mid C\mid\sigma(q^{a_q})\mid2N$. Since $r$ is odd, $r\mid N$, contradicting the definition of $P$.
\end{proof}

\begin{lemma}\label{lem:valuation}
For $q\in S\setminus\{P\}$,
\[
 P\mid\sigma(q^{a_q})\quad\Longleftrightarrow\quad d_q\mid a_q+1.
\]
When these equivalent conditions hold,
\begin{equation}\label{eq:valuation}
 w_q=v_P(q^{d_q}-1)=v_P(q^{P-1}-1).
\end{equation}
Moreover,
\begin{equation}\label{eq:balance}
 a_P=\sum_{q\in S\setminus\{P\}}w_q.
\end{equation}
\end{lemma}
\begin{proof}
Since $q<P$, the prime $P$ does not divide $q-1$. The identity
\[
 \sigma(q^{a_q})=\frac{q^{a_q+1}-1}{q-1}
\]
therefore gives the equivalence. If $a_q+1=d_qm$, Lemma~\ref{lem:noexponent} implies $P\nmid m$, while
\[
 \frac{q^{d_qm}-1}{q^{d_q}-1}
 =1+q^{d_q}+\cdots+q^{d_q(m-1)}\equiv m\not\equiv0\pmod P.
\]
This proves the first equality in~\eqref{eq:valuation}. The second follows by the same argument with $m=(P-1)/d_q$.

Finally, $\sigma(P^{a_P})\equiv1\pmod P$. Taking $P$-adic valuations in $\prod_{q\in S}\sigma(q^{a_q})=2N$ proves~\eqref{eq:balance}.
\end{proof}

\section{The contributions of orders three and six}

\begin{lemma}\label{lem:unit}
If $q<P$ contributes to $P$ and $d_q\in\{3,6\}$, then $w_q=1$.
\end{lemma}
\begin{proof}
If $d_q=3$, then $P\mid q^2+q+1$. The primes $q$ and $P$ are odd, so $q\le P-2$, and
\[
 0<q^2+q+1\le P^2-3P+3<P^2.
\]
It follows that $v_P(q^2+q+1)=1$. As $P\nmid q-1$, Lemma~\ref{lem:valuation} gives $w_q=1$.

If $d_q=6$, then $P\mid q^2-q+1$. In the factorization
\[
 q^6-1=(q-1)(q+1)(q^2+q+1)(q^2-q+1),
\]
the first three factors are not divisible by $P$: divisibility of any of them would give order $1$, $2$, or $3$. Since $0<q^2-q+1<P^2$, equation~\eqref{eq:valuation} again gives $w_q=1$.
\end{proof}

\begin{lemma}\label{lem:incompatible}
There do not exist odd primes $u,v<P$ satisfying
\[
 \ord_P(u)=3,\qquad \ord_P(v)=6.
\]
\end{lemma}
\begin{proof}
Such primes would satisfy $P\mid u^2+u+1$ and $P\mid v^2-v+1$. Subtracting yields
\[
 P\mid(u+v)(u-v+1).
\]
The positive integer $u+v$ is even and less than $2P$, so it is not divisible by the odd prime $P$. The integer $u-v+1$ is odd and has absolute value less than $P$, so it too is not divisible by $P$. This is a contradiction.
\end{proof}

\begin{lemma}\label{lem:count}
At most two primes contribute to $P$ through order $3$, and at most one contributes through order $6$. If $3\nmid N$, at most one contributes through order $3$.
\end{lemma}
\begin{proof}
Primes of order $3$ are roots of $x^2+x+1$ modulo $P$. This polynomial has at most two roots, and distinct primes below $P$ represent distinct residues.

If a prime $q$ contributes through order $6$, then $6\mid a_q+1$, so $a_q$ is odd. For an odd base $q$, the sum $\sigma(q^{a_q})$ is then even. Two distinct such components would make $4\mid\sigma(N)=2N$, contrary to the oddness of $N$. Thus there is at most one.

For the last assertion, suppose $3\nmid N$ and two distinct primes $u,v<P$ contribute through order $3$. The congruences for the two roots give
\[
 P\mid(u-v)(u+v+1).
\]
Since $u\ne v$ and $0<u+v+1<2P$, this forces
\begin{equation}\label{eq:pair}
 u+v+1=P.
\end{equation}
Order $3$ implies $P\equiv1\pmod3$. Neither $u$ nor $v$ equals $3$, and~\eqref{eq:pair} shows that their residues modulo $3$ are $1$ and $2$. Suppose $u\equiv1\pmod3$. Because $3\mid a_u+1$, one has
\[
 \sigma(u^{a_u})\equiv a_u+1\equiv0\pmod3.
\]
As $\sigma(u^{a_u})\mid2N$, this contradicts $3\nmid N$.
\end{proof}

\begin{proof}[Proof of Theorem~\ref{thm:main}]
By Lemma~\ref{lem:incompatible}, the contributing primes cannot include both an order-$3$ prime and an order-$6$ prime. Lemmas~\ref{lem:unit} and~\ref{lem:count} therefore give~\eqref{eq:lowbound}. Subtracting that bound from~\eqref{eq:balance} proves~\eqref{eq:main}.
\end{proof}

\section{Consequences}

\begin{corollary}\label{cor:fourth}
If $P^4\mid N$, then at least one of the following holds:
\begin{itemize}
\item There are distinct primes $u,v\in H$ such that
\[
 P\mid\sigma(u^{a_u}),\qquad P\mid\sigma(v^{a_v}).
\]
\item There is a prime $q\in H$ such that
\[
 P\mid\sigma(q^{a_q}),\qquad q^{P-1}\equiv1\pmod{P^2}.
\]
\end{itemize}
\end{corollary}
\begin{proof}
Theorem~\ref{thm:main} gives $\sum_{q\in H}w_q\ge2$. Either at least two summands are positive, or a single summand is at least two. In the latter case, equation~\eqref{eq:valuation} gives $P^2\mid q^{P-1}-1$.
\end{proof}

\begin{corollary}\label{cor:exception}
Suppose $P\nmid\sigma(q^{a_q})$ for every $q\in H$. Then either $a_P=1$, or $a_P=2$, $3\mid N$, and there are distinct primes $u,v<P$ dividing $N$ such that
\[
 d_u=d_v=3,\qquad 3\mid a_u+1,\qquad 3\mid a_v+1,
\]
\[
 u+v+1=P,\qquad P\mid u^2+u+1,\qquad P\mid v^2+v+1.
\]
\end{corollary}
\begin{proof}
Theorem~\ref{thm:main} gives $a_P\le2$. If $a_P=2$, the improved bound implies $3\mid N$. The two required contributions cannot come from order $6$, which has at most one contributing prime. They therefore come from two distinct primes of order $3$. Lemma~\ref{lem:valuation} gives the shifted-exponent divisibilities, and the calculation in~\eqref{eq:pair} gives their sum.
\end{proof}

\begin{corollary}\label{cor:exponents}
If $q\in H$ contributes to $P$, then $a_q\ge4$. If $a_q$ is odd, then $a_q\ge9$; in particular, this applies when $q$ is the Euler prime, the prime occurring to an odd exponent.
\end{corollary}
\begin{proof}
An odd prime $q<P$ cannot have order $1$ or $2$ modulo $P$: these would require $q=1$ or $q=P-1$, respectively. Also $4\nmid a_q+1$. Indeed, if $4\mid a_q+1$, then
\[
 (q+1)(q^2+1)=\sigma(q^3)\mid\sigma(q^{a_q}),
\]
and both factors on the left are even, contradicting $\sigma(q^{a_q})\mid2N$.

Since $d_q\mid a_q+1$, exponents $1$ and $2$ would give orders $2$ and $3$, while exponent $3$ is excluded by the preceding argument. Thus $a_q\ge4$. If $a_q$ is odd and less than $9$, the only remaining possibilities are $5$ and $7$. Exponent $5$ allows only orders $1,2,3,6$, and exponent $7$ makes $4\mid a_q+1$. Both are impossible.
\end{proof}

\section{The 3-adic and Euler-prime refinements}\label{sec:three}

Set
\[
 W=\sum_{q\in H}w_q,\qquad L=a_P-W,\qquad
 b=v_3(N),\qquad t=v_3(a_P+1).
\]
There is exactly one prime divisor $e$ of $N$ for which $a_e$ is odd. Indeed,
$\sigma(q^{a_q})$ is even precisely when $a_q$ is odd, and the product of these
components is $2N$, which has $2$-adic valuation one. We call $e$ the Euler prime.
Define
\[
 \delta=\begin{cases}
 v_3(e+1)+v_3(a_e+1),& e\equiv2\pmod3,\\
 0,&\text{otherwise}.
 \end{cases}
\]
The standard lifting-the-exponent identities give
\begin{align}
 v_3(\sigma(q^{a_q}))&=v_3(a_q+1)
       &&\text{if }q\equiv1\pmod3,\label{eq:threeone}\\
 v_3(\sigma(e^{a_e}))&=\delta
       &&\text{if }e\equiv2\pmod3.\label{eq:threetwo}
\end{align}
These are instances of the classical component valuation formulas; see
Voight~\cite{Voight}. Also
\begin{equation}\label{eq:threebalance}
 \sum_{q\in S}v_3(\sigma(q^{a_q}))=b.
\end{equation}

\begin{theorem}\label{thm:three}
For $s=t+\delta$, define
\[
 A(b,s)=\begin{cases}
 0,& b<s,\\
 1,& b=s,\\
 2,& b>s.
 \end{cases}
\]
Then
\[
 \boxed{W\ge a_P-A(b,t+\delta).}
\]
In particular, $W\ge a_P-A(b,t)$.
\end{theorem}
\begin{proof}
The original bound gives $0\le L\le2$. If $L>0$, there is a contributing prime
of order $3$ or $6$ modulo $P$, so $P\equiv1\pmod3$. The component at $P$
then contributes $t$ to~\eqref{eq:threebalance}. If $e\equiv2\pmod3$, the Euler
prime differs from $P$ and contributes an additional $\delta$. Thus
\[
 L>0\quad\Longrightarrow\quad b\ge t+\delta.
\]
If $L=2$, Lemmas~\ref{lem:incompatible} and~\ref{lem:count} give two distinct
contributing order-$3$ primes $u,v<P$ with $u+v+1=P$. Neither is $3$: otherwise
the other would also be $0$ modulo $3$, hence also $3$. Their residues modulo $3$
are therefore $1$ and $2$. Label them so that $u\equiv1\pmod3$.
Since $3\mid a_u+1$, equation~\eqref{eq:threeone} supplies at least one further
unit of $3$-adic valuation. This component differs from $P$ and, when
$e\equiv2\pmod3$, also from $e$. Consequently
\[
 L=2\quad\Longrightarrow\quad b\ge t+\delta+1.
\]
These two implications prove the displayed piecewise bound. Omitting the
nonnegative Euler-prime contribution gives the bound with $t$ alone.
\end{proof}

For example, if $a_3=2$ and $a_P=8$, then $b=t=2$, so $W\ge7$.
If also $e\equiv2\pmod3$, then $\delta\ge1$ and $W=8$.
If $3\nmid N$ and $a_P=2$, then $b=0<t=1$, so $W=2$; the argument of
Corollary~\ref{cor:fourth} then gives two contributing primes in $H$ or a
contributing prime satisfying the congruence modulo $P^2$.

\section{Distinct fourth-power prime divisors}\label{sec:fourthcount}

Define
\[
 K=\#\{r\in S:r<P,\ r\equiv1\pmod3,\ a_r\ge4\}.
\]
The following condition couples the preceding valuation restriction to the
number of distinct prime divisors occurring at least to the fourth power.

\begin{theorem}\label{thm:fourthcount}
If $P\equiv1\pmod3$, then
\begin{equation}\label{eq:fourthcount}
 \boxed{K+b\ge2t+\delta.}
\end{equation}
In particular, $K+b\ge2t$. If $b=t+\delta$, then $K\ge t$.
\end{theorem}
\begin{proof}
For $0\le j<t$, put
\[
 x_j=P^{3^j},\qquad
 C_j=\frac{x_j^2+x_j+1}{3}=\frac{\Phi_{3^{j+1}}(P)}{3}.
\]
We first establish four properties of these integers.

Since $x_j\equiv1\pmod3$, write $x_j=3k+1$. Then
\[
 x_j^2+x_j+1=3(3k^2+3k+1),\qquad C_j\equiv1\pmod3.
\]
In particular, $3\nmid C_j$. Also $x_j\ge P>1$ and
$x_j^2+x_j+1>3x_j$, so $C_j>P$.

Next, $3^{j+1}\mid a_P+1$, and the geometric-series factorization gives
\[
 x_j^2+x_j+1\mid\sigma(P^{a_P}).
\]
Thus $C_j\mid2N$. The numerator and $C_j$ are odd, so $C_j\mid N$.

If a prime $r$ divides $C_j$, then $r\ne3$ and
$x_j^2+x_j+1\equiv0\pmod r$. Hence $x_j^3\equiv1\pmod r$, whereas
$x_j\not\equiv1\pmod r$, since the latter would force $r=3$.
It follows that $\ord_r(x_j)=3$ and $r\equiv1\pmod3$.
Moreover, $P$ does not divide $C_j$, so maximality of $P$ gives $r<P$.

Finally, suppose $i<j$ and a prime $r$ divides both $C_i$ and $C_j$.
Divisibility of $C_i$ implies $P^{3^{i+1}}\equiv1\pmod r$, and hence
$P^{3^j}\equiv1\pmod r$. Divisibility of $C_j$ would then give
$0\equiv1+1+1\pmod r$, contrary to $r\ne3$. The integers $C_j$
are therefore pairwise coprime.

Every $C_j$ has a prime divisor $r_j$ for which $a_{r_j}\ge2$.
Otherwise all its prime divisors would occur in $N$ to exponent one. Each
would have to be the unique prime $e$ with odd exponent, and $C_j\mid N$
would force $C_j\mid e$. This contradicts $C_j>P\ge e$.
Choose one such $r_j$ for each $j$. The coprimality just proved makes these
$t$ primes distinct, and all satisfy $r_j<P$ and $r_j\equiv1\pmod3$.

Exponent three is impossible by the argument in Corollary~\ref{cor:exponents}.
Thus each selected prime either has exponent at least four or has exponent
exactly two. Let $R$ be the number with exponent two. By~\eqref{eq:threeone},
each of those contributes exactly one to~\eqref{eq:threebalance}.
The component at $P$ contributes $t$. When $e\equiv2\pmod3$, its component
is disjoint from all the selected primes and contributes $\delta$.
Consequently
\[
 R+t+\delta\le b.
\]
At least $t-R$ of the selected primes are counted by $K$. Therefore
\[
 K\ge t-R\ge2t+\delta-b,
\]
which proves~\eqref{eq:fourthcount} and both stated consequences.
\end{proof}

\begin{corollary}\label{cor:eightfourth}
If $a_3=2$ and $a_P=8$, then
\[
 \boxed{W=8\quad\text{or}\quad(W=7\ \text{and}\ K\ge2).}
\]
\end{corollary}
\begin{proof}
Theorem~\ref{thm:three} gives $7\le W\le8$. If $W=7$, then $L=1>0$,
so $P\equiv1\pmod3$. Since $b=t=2$, Theorem~\ref{thm:fourthcount} gives
$K+2\ge4$, hence $K\ge2$.
\end{proof}

This adds a restriction to $W\ge7$: a candidate exponent pattern with
$W=7$ and fewer than two such fourth-power prime divisors is excluded.
The two primes are distinct even if the contribution measured by $W$ is
concentrated at a single prime. They are not asserted to be members of $H$.
The cyclotomic and valuation ingredients are classical; the formal proof
establishes the stated combination, not its publication priority.

\section{Upper bounds for the contribution outside orders three and six}\label{sec:upper}

The valuation balance immediately gives $W\le a_P$. We now obtain upper bounds
which do not use $a_P$. Let
\[
 T=\{q\in H:w_q>0\},\qquad M=|T|,
\]
and let $\varphi$ denote Euler's totient function. By Lemma~\ref{lem:valuation},
membership in $T$ is equivalent to $q\in H$ and $d_q\mid a_q+1$.

\begin{theorem}\label{thm:upper}
For each $q\in T$,
\begin{equation}\label{eq:localupper}
 w_q=v_P\bigl(\Phi_{d_q}(q)\bigr)
 \le\left\lfloor\log_P\Phi_{d_q}(q)\right\rfloor
 \le\varphi(d_q)-1.
\end{equation}
Consequently,
\begin{equation}\label{eq:sumupper}
 \boxed{W\le
 \sum_{q\in T}\left\lfloor\log_P\Phi_{d_q}(q)\right\rfloor
 \le\sum_{q\in T}\bigl(\varphi(d_q)-1\bigr).}
\end{equation}
\end{theorem}
\begin{proof}
Fix $q\in T$ and put $d=d_q$. Lemma~\ref{lem:valuation} gives
$w_q=v_P(q^d-1)$. In the factorization
\[
 q^d-1=\prod_{r\mid d}\Phi_r(q),
\]
only the factor with $r=d$ can be divisible by $P$. Indeed, if
$P\mid\Phi_r(q)$, then $P\mid q^r-1$, so $d\mid r$. Since also $r\mid d$,
we must have $r=d$. This proves the equality in~\eqref{eq:localupper}.

The order $d$ is neither $1$ nor $2$, by the argument in
Corollary~\ref{cor:exponents}. The cyclotomic value is positive and
\[
 \Phi_d(q)=\prod_{\substack{\zeta\text{ primitive}\\d\text{-th root of unity}}}
 |q-\zeta|\le(q+1)^{\varphi(d)}<P^{\varphi(d)}.
\]
For the strict inequality, use $q+1<P$, since $q<P$ are odd primes.
As $P^{w_q}\mid\Phi_d(q)$, we have
\[
 P^{w_q}\le\Phi_d(q)<P^{\varphi(d)}.
\]
Taking integer logarithms proves both inequalities in~\eqref{eq:localupper}.
The primes in $H\setminus T$ have zero contribution, so summation proves
\eqref{eq:sumupper}.
\end{proof}

\begin{corollary}\label{cor:oddpartupper}
Write $P-1=2^s m$, with $m$ odd. Then
\[
 \boxed{W\le M\bigl(\varphi(m)-1\bigr).}
\]
\end{corollary}
\begin{proof}
For $q\in T$, we have $d_q\mid P-1$ and $d_q\mid a_q+1$.
The proof of Corollary~\ref{cor:exponents} gives $4\nmid a_q+1$, hence
$4\nmid d_q$. It follows that $d_q=u$ or $d_q=2u$, where $u$ is an odd
divisor of $m$. In either case $\varphi(d_q)=\varphi(u)$. The prime-power
formula for the totient gives $\varphi(u)\mid\varphi(m)$, so
$\varphi(d_q)\le\varphi(m)$. Apply Theorem~\ref{thm:upper} to each of the
$M$ contributing primes.
\end{proof}

\begin{corollary}\label{cor:exponentupper}
With $c(N)$ as in Theorem~\ref{thm:main},
\[
 a_P\le c(N)+\sum_{q\in T}\bigl(\varphi(d_q)-1\bigr)
 \le c(N)+M\bigl(\varphi(m)-1\bigr).
\]
\end{corollary}
\begin{proof}
Combine the lower bound $a_P-c(N)\le W$ with
Theorem~\ref{thm:upper} and the proof of Corollary~\ref{cor:oddpartupper}.
\end{proof}

For instance, if $P-1=2^s\cdot5$, every contributing order outside $3$ and $6$
is $5$ or $10$. Each contribution is at most three, and therefore
\[
 W\le3M,\qquad a_P\le3M+c(N).
\]
These bounds can be combined with $W\le a_P$ by taking the smaller upper
bound. Their additional content is a restriction in terms of the smaller
prime divisors and their orders, independently of the exponent at $P$.

\begin{corollary}\label{cor:nowieferich}
If $q^{P-1}\not\equiv1\pmod{P^2}$ for every $q\in T$, then $W=M$.
\end{corollary}
\begin{proof}
By Lemma~\ref{lem:valuation}, every $q\in T$ satisfies
$w_q=v_P(q^{P-1}-1)\ge1$. The stated condition makes this valuation exactly
one. Summing gives the result.
\end{proof}

The last hypothesis cannot be discarded merely because $q<P$.
For example, $\ord_{11}(3)=5$ and $\Phi_5(3)=121=11^2$.
This is an illustration of the local arithmetic, not an odd perfect number.
The estimates above use classical cyclotomic facts and are not claimed as
new results in the literature.

\begin{theorem}[Combined inequality]\label{thm:combined}
Let $N$ be an odd perfect number, so that
$\sigma(N)=\sum_{d\mid N}d=2N$, and write
\[
 N=\prod_{q\in S}q^{a_q},\qquad
 S=\{q:q\text{ is prime and }q\mid N\},\qquad P=\max S.
\]
For $q\in S\setminus\{P\}$, define
$d_q=\ord_P(q)$, the multiplicative order of $q$ modulo $P$, and
$w_q=v_P(\sigma(q^{a_q}))$, where $v_r(x)$ denotes the exponent of the prime
$r$ in $x$. Set
\[
 H=\{q\in S:q<P,\ d_q\notin\{3,6\}\},\quad
 W=\sum_{q\in H}w_q,\quad
 T=\{q\in H:w_q>0\},\quad M=|T|.
\]
Let $e\in S$ be the unique prime with $a_e$ odd, and define
\[
 b=v_3(N),\qquad t=v_3(a_P+1),\qquad
 \delta=\begin{cases}
 v_3(e+1)+v_3(a_e+1),&e\equiv2\pmod3,\\
 0,&\text{otherwise},
 \end{cases}
\]
\[
 A(b,u)=\begin{cases}
 0,&b<u,\\ 1,&b=u,\\ 2,&b>u.
 \end{cases}
\]
Write $P-1=2^s m$ with $m$ odd. Let $\Phi_d$ denote the $d$-th cyclotomic
polynomial and $\varphi(d)$ Euler's totient function. Then
\begin{equation}\label{eq:combined}
\boxed{\begin{aligned}
 a_P-A(b,t+\delta)&\le W\\
 &\le\min\!\left\{a_P,\sum_{q\in T}
   \left\lfloor\log_P\Phi_{d_q}(q)\right\rfloor\right\}\\
 &\le\min\!\left\{a_P,\sum_{q\in T}
   \bigl(\varphi(d_q)-1\bigr)\right\}\\
 &\le\min\!\left\{a_P,M\bigl(\varphi(m)-1\bigr)\right\}.
\end{aligned}}
\end{equation}
In particular, $a_P\le A(b,t+\delta)+M(\varphi(m)-1)$.
\end{theorem}
\begin{proof}
Theorem~\ref{thm:three} gives the lower bound. The valuation identity
$a_P=\sum_{q\in S\setminus\{P\}}w_q$ gives $W\le a_P$ because all valuations
are nonnegative. Theorem~\ref{thm:upper} gives the logarithmic and totient
upper bounds. Corollary~\ref{cor:oddpartupper} bounds each totient summand by
$\varphi(m)-1$, proving the final inequality. Comparing the endpoints yields
the stated consequence for $a_P$.
\end{proof}

\section{Formal verification}\label{sec:lean}

The accompanying files use Lean~4.19.0 and mathlib revision
\begin{center}\small\texttt{c44e0c8ee63ca166450922a373c7409c5d26b00b}.\end{center}
The formal statements take \texttt{Nat.Perfect N}, \texttt{Odd N}, membership of $P$ in \texttt{N.primeFactors}, and maximality of $P$ among those factors as hypotheses. The divisor sums, multiplicative orders, factorization exponents, and valuations are the corresponding arithmetic objects; no valuation balance or order restriction is supplied as an extra hypothesis.

The definition \texttt{component N q} is the geometric sum
\[
 \sum_{i=0}^{N.\mathrm{factorization}(q)}q^i.
\]
The lemma \texttt{component\_sigma} identifies it with mathlib's divisor-sum function on the prime power. The definition \texttt{ord P q} is \texttt{orderOf (q : ZMod P)}, and \texttt{high N P} is the set $H$ in Section~1.

\begin{center}
\small
\begin{tabular}{@{}p{0.27\textwidth}p{0.67\textwidth}@{}}
\toprule
Result & Lean declaration \\
\midrule
General bound & \texttt{largestPrime\_order\_bound} \\[3pt]
Bound if $3\nmid N$ & \texttt{largestPrime\_order\_bound\_without\_three} \\[3pt]
Corollary~\ref{cor:fourth} & \texttt{fourth\_power\_alternative} \\[3pt]
Corollary~\ref{cor:exception} & \texttt{exceptional\_pair} \\[3pt]
Corollary~\ref{cor:exponents} & \texttt{high\_contributor\_exponents} \\
\bottomrule
\end{tabular}
\end{center}

The bounds in~\eqref{eq:lowbound} are also stated explicitly as
\begin{center}\small
\texttt{low\_order\_contribution\_bound}\\
\texttt{low\_order\_contribution\_bound\_without\_three}.
\end{center}
Their sums run over contributing primes; the omitted terms have valuation zero.

The additional proof modules implement Sections~\ref{sec:three}
and~\ref{sec:fourthcount}. They derive the divisibility, size, residue,
disjointness, and counting arguments from the same odd-perfect-number
hypotheses. The Euler prime's existence and uniqueness are also proved
in the project. The README identifies each module and its role.
The main declarations are:
\begin{center}
\small
\begin{tabular}{@{}p{0.32\textwidth}p{0.62\textwidth}@{}}
\toprule
Result & Lean declaration \\
\midrule
Theorem~\ref{thm:three}, without $\delta$ & \texttt{three\_adic\_refinement} \\[3pt]
Theorem~\ref{thm:three}, with $\delta$ & \texttt{euler\_three\_adic\_refinement} \\[3pt]
Theorem~\ref{thm:fourthcount}, with $\delta$ & \texttt{fourth\_power\_budget\_euler} \\[3pt]
Theorem~\ref{thm:fourthcount}, without $\delta$ & \texttt{fourth\_power\_budget} \\[3pt]
Case $b=t+\delta$ & \texttt{saturated\_fourth\_power\_count} \\[3pt]
Corollary~\ref{cor:eightfourth} & \texttt{eighth\_power\_fourth\_power\_dichotomy} \\
\bottomrule
\end{tabular}
\end{center}

Natural-number subtraction in the two main Lean inequalities replaces a negative lower bound by zero. This is equivalent to the displayed mathematical inequality, since the sum of valuations is nonnegative. The hypothesis $P^4\mid N$ in Corollary~\ref{cor:fourth} is represented by $4\le N.\mathrm{factorization}(P)$, and its congruence by the equivalent divisibility $P^2\mid q^{P-1}-1$. The formal exceptional-pair statement uses membership in the set of contributing order-$3$ primes; the polynomial and shifted-exponent divisibilities follow from the included lemmas.

The verification entry point is \texttt{Audit.lean}. It checks the named declarations and prints their axiom dependencies. The source contains no \texttt{sorry}, \texttt{admit}, additional axioms, or use of \texttt{native\_decide}. The audit reports only Lean's standard axioms \texttt{propext}, \texttt{Classical.choice}, and \texttt{Quot.sound}. Compilation output and reproduction instructions are included in the archive. Formal verification establishes the stated implications; it does not establish publication priority.

The upper bounds in Section~\ref{sec:upper} are formalized in
\texttt{CyclotomicUpper.lean}.

The set \texttt{highContributors N P} is $T$;
it is defined by filtering $H$ for divisibility of the component by $P$.
The definition \texttt{cycloValue q d} is the absolute value of the integer
$\Phi_d(q)$, whose positivity for $q>1$ is proved using mathlib.
Integer logarithms are represented by \texttt{Nat.log}.

\begin{center}
\small
\begin{tabular}{@{}p{0.32\textwidth}p{0.62\textwidth}@{}}
\toprule
Result & Lean declaration \\
\midrule
Component valuation identity & \texttt{component\_val\_eq\_cycloValue} \\[3pt]
Logarithmic upper bound & \texttt{highWeight\_log\_upper} \\[3pt]
Comparison of bounds & \texttt{logarithmic\_bound\_le\_totient\_bound} \\[3pt]
Totient upper bound & \texttt{highWeight\_totient\_upper} \\[3pt]
Corollary~\ref{cor:oddpartupper} & \texttt{highWeight\_odd\_part\_upper} \\[3pt]
Corollary~\ref{cor:exponentupper}, first bound & \texttt{largest\_exponent\_totient\_upper} \\[3pt]
Corollary~\ref{cor:nowieferich} & \texttt{highWeight\_eq\_card\_of\_no\_wieferich} \\
Theorem~\ref{thm:combined} & \texttt{combined\_inequality} \\[3pt]
Its exponent consequence & \texttt{combined\_exponent\_constraint} \\
\bottomrule
\end{tabular}
\end{center}

\begin{thebibliography}{9}
\bibitem{ClaytonHansen}
G.~Clayton and C.~Hansen,
\emph{On inequalities involving counts of the prime factors of an odd perfect number},
Integers \textbf{23} (2023), Paper A79.
\href{https://math.colgate.edu/~integers/x79/x79.pdf}{Full text}.

\bibitem{Voight}
J.~Voight,
\emph{On the nonexistence of odd perfect numbers},
MASS Selecta, American Mathematical Society, 2003, pp.~293--300.
\href{https://jvoight.github.io/articles/opn-mass-rev-060211.pdf}{Author's version}.
\end{thebibliography}
\end{document}
